\chapter{Phonemes}

\section{Vowels}

\subsection{Monophthongs}

Table~\vref{tbl:vowel-inventory} contains Interturkic's eight vowel phonemes,
arranged according to their contrasts in height, backness, and rounding.

\begin{table}
  \centering
  \caption{Vowel inventory}
  \label{tbl:vowel-inventory}
  \begin{tabular}{lllll}
    \toprule
    &\multicolumn{2}{l}{Front}&\multicolumn{2}{l}{Back}\\
    \cmidrule(lr){2-3}\cmidrule(lr){4-5}
    &unrounded&rounded&unrounded&rounded\\
    \midrule
    High&\phonem{i}&\phonem{y}&\phonem{ɯ}&\phonem{u}\\
    Low&\phonem{ɛ}&\phonem{œ}&\phonem{ɑ}&\phonem{ɔ}\\
    \bottomrule
  \end{tabular}
\end{table}

In native vocabulary, \phonem{ɔ, œ} occur only in the first syllables of roots.
All of the other vowels are unrestricted in their distribution, and may occur in
any position.

\subsection{Diphthongs}

Interturkic has small set of diphthongs, although they are not phonemic in
their own right.
They are formed from and analysed as monophthongs followed by glides~(\vref
{ex:diphthong-yes}).
Note that such sequences are not considered diphthongs if split across syllable
boundaries~(\vref{ex:diphthong-no}).

\pex
  \a\label{ex:diphthong-yes}%
    \vtop{\halign{\phonem{#}\hfil&\:\transl{#}\hfil\cr
      ɑj&moon, month\cr
      tɑw&mountain\cr
    }}
  \a\label{ex:diphthong-no}%
    \vtop{\halign{\phonem{#}\hfil&\:\transl{#}\hfil\cr
      bɔ.jun&neck\cr
      ɑ.wul&village\cr
    }}
\xe

Diphthongs formed with \phonem{w} are treated as if they are rounded for the
purposes of vowel harmony.

\section{Consonants}

\begin{table}
  \centering
  \caption{Consonant inventory}
  \begin{tabular}{lll}
    \toprule
    &Unvoiced&Voiced\\
    \midrule
    Nasal&&\phonem{m n ŋ}\\
    Stop&\phonem{p t k q ʔ}&\phonem{b d g}\\
    Affricate&\phonem{t͡ʃ}&\phonem{d͡ʒ}\\
    Fricative&\phonem{f s ʃ x h}&\phonem{v z ʒ ɣ}\\
    Liquid&&\phonem{ɾ l}\\
    Glide&&\phonem{w j}\\
    \bottomrule
  \end{tabular}
\end{table}
